%%%%%%%%%%%%%%
%  COSTANTI  %
%%%%%%%%%%%%%%
%COMANDI DA UTILIZZARE PER TUTTI I DOCUMENTI
% In questa prima parte vanno definite le 'costanti' utilizzate da due o più documenti.

%COMANDI TABELLE
\newcommand{\rowcolorhead}{\rowcolor[HTML]{00ACC1}} %intestazione 
\newcommand{\rowcolorlight}{\rowcolor[HTML]{E0f7fa}} %righe chiare/dispari
\newcommand{\rowcolordark}{\rowcolor[HTML]{80DEEA}} %righe scure/pari
\newcommand{\colorhead}{\color[HTML]{FFFFFF}} %testo intestazione
\newcommand{\colorbody}{\color[HTML]{000000}} %testo righe


%%%%%%%%%%%%%%
%  FUNZIONI  %
%%%%%%%%%%%%%%

% Serve a dare la giusta formattazione al codice inline
\newcommand{\code}[1]{\flextt{\texttt{#1}}}

% Genera automaticamente la pagina di copertina
\newcommand{\makeFrontPage}{
  % Declare new geometry for the title page only.
  \newgeometry{top=4cm}
  
  \begin{titlepage}
  \begin{center}

\begin{center}
  \centerline{\includegraphics[scale=0.24]{StyleLatex/logo_home.png}}
\end{center}
  
  \vspace{1cm}

  \begin{Huge}
  \textbf{\DocTitolo{}} \\
  \end{Huge}

  \vspace{9pt}  
  
  \begin{tabular}{l|l}
  	\textbf{Nome Gruppo} & I Tesori di Squitty \\
  	\textbf{Componenti} & Alessio Gobbo 1125947\\
  	& Simone Ballarin 1122286\\
  	& Riccardo Dario\\
  	& Gerta Lleshi\\
  	\textbf{Dati login utente} & username: kevinsilvstr \\
  	& password: Kevin940 \\
  	\textbf{Dati login admin} & username: admin \\
  	& password: Admin420 \\
  	\textbf{Username consegna} & ksilvest\\ 
  	\hline
  \end{tabular}

  \vspace{30pt}
  
\begin{center}
	\centerline{Indirizzo sito: 
	http://tecweb2016.studenti.math.unipd.it/ksilvest/WineNot/}
\end{center}

  \end{center}
  \end{titlepage}
  
  % Ends the declared geometry for the titlepage
  \restoregeometry
}
%\documentclass[a4paper, oneside, openany]{book}
\documentclass[10pt,a4paper]{extarticle}

%**************************************************************
% Importazione package
%************************************************************** 
\usepackage{graphbox}
% permette di modificare i margini
%\usepackage[top=3.1cm, bottom=3.1cm, left=2.2cm, right=2.2cm]{geometry}
%\usepackage[a4paper]{../../common/style/geometry} %TODO
\usepackage[a4paper, top=2cm, bottom=2.8cm]{geometry}
\usepackage{booktabs} %lista item nelle tabelle
% specifica con quale codifica bisogna leggere i file
\usepackage[utf8]{inputenc}

%indice con i puntini
\usepackage{tocloft}
\renewcommand\cftsecleader{\cftdotfill{\cftdotsep}}



% necessario per risolvere il problema del carattere invisibile per l'a capo
\DeclareUnicodeCharacter{00A0}{ } %?

% per scrivere in italiano e in inglese;
% l'ultima lingua (l'italiano) risulta predefinita
\usepackage[english, italian]{babel}

% imposta lo stile italiano per i paragrafi
\usepackage{parskip}

% fornice elenchi numerati in ordine inverso
\usepackage{etaremune}

\usepackage{caption}

% comandi per l'appendice
\usepackage{appendix}
%\renewcommand\appendixtocname{Appendici} %We really need it?

% import euro symbol
\usepackage{eurosym}

% numera anche i sottoparagrafi
\setcounter{secnumdepth}{5}

% elenca anche i sottoparagrafi nell'indice
\setcounter{tocdepth}{5}

% permetti di definire dei colori
\usepackage[usenames,dvipsnames]{color}

% permette di usare il comando "paragraph" come subsubsubsection!
\usepackage{titlesec}

% permette di inserire le immagini/tabelle esattamente dove viene usato il
% comando \begin{figure}[H] ... \end{figure}
% evitando che venga spostato in automatico
\usepackage{float}

% permette l'inserimento di url e di altri tipi di collegamento
\usepackage[colorlinks=true]{hyperref}

\hypersetup{
    colorlinks=true, % false: boxed links; true: colored links
    citecolor=black,
    filecolor=black,
    linkcolor=black, % color of internal links
    urlcolor=Maroon  % color of external links
}

% permette al comando \url{...} di andare a capo a metà di un link
\usepackage{breakurl}

% immagini
\usepackage{graphicx}

% permette di riferirsi all'ultima pagina con \pageref{LastPage}
\usepackage{lastpage}

% tabelle su più pagine
\usepackage{longtable}

% per avere dei comandi in più da poter usare sulle tabelle
\usepackage{booktabs}

% tabelle con il campo X per riempire lo spazio rimanente sulla riga
%\usepackage{tabularx}

% TABELLE 
% tabelle su più pagine
\usepackage{longtable}

% per avere dei comandi in più da poter usare sulle tabelle
\usepackage{booktabs}

% multirow per tabelle
\usepackage{multirow}

% definisci un nuovo tipo di colonna C che permette di andare a capo con \newline
% e centrata
\usepackage{array}
\usepackage{ragged2e}
%\newcolumntype{P}[1]{>{\RaggedRight\hspace{0pt}}p{#1}}
\newcolumntype{C}[1]{>{\centering\let\newline\\\arraybackslash\hspace{0pt}}m{#1}}

% colore di sfondo per le celle
\usepackage[usenames,dvipsnames,svgnames,table]{xcolor}

% tabelle con il campo X per riempire lo spazio rimanente sulla riga
\usepackage{tabularx}


% personalizza l'intestazione e piè di pagina
\usepackage{fancyhdr}

% permette di inserire caratteri speciali
\usepackage{textcomp}

% permette di aggiustare i margini e centrare tabelle e figure
\usepackage{changepage}

%Permette di includere i grafici a barre
%IMPORTANTE: deve essere caricato prima di /pgfgantt altrimenti causa conflitto
%\usepackage{pgfplots}

% permette di includere i diagrammi Gantt
% su Ubuntu non si può installare il pacchetto, deve essere in modello/
%\usepackage{../../modello/pgfgantt}

% permette di includere i grafici a torta
%\usepackage{../../modello/pgf-pie}

% necessario per pgf-pie
\usepackage{tikz}

% permette i path delle immagini con gli spazi
\usepackage{grffile}

% ruota le immagini
\usepackage{rotating}

% permetti di calcolare le larghezze facendo calcoli
\usepackage{calc}


\fancypagestyle{plain}{
	% cancella tutti i campi di intestazione e piè di pagina
	\fancyhf{}

	\lfoot{
		\DocTitolo{} \ - \textit{\DocRedazione{} \ -  \DocAnno{}}
	}
	\rfoot{Pagina \thepage{} di \pageref{LastPage}}

	% Visualizza una linea orizzontale in cima e in fondo alla pagina
	\renewcommand{\headrulewidth}{0pt}	
	\renewcommand{\footrulewidth}{0.3pt}
}
\pagestyle{plain}

% allarga l'header a tutta la pagina
%\fancyhfoffset[L]{\oddsidemargin + \hoffset + 1in}
%\fancyhfoffset[R]{\evensidemargin + \marginparwidth - \marginparsep}

% Per inserire del codice sorgente formattato

\usepackage{listings}

\lstset{
  extendedchars=true,          % lets you use non-ASCII characters
  inputencoding=utf8,   % converte i caratteri utf8 in latin1, richiede 
  %\usepackage{listingsutf8} anzichè listings
  basicstyle=\ttfamily,        % the size of the fonts that are used for the 
  %code
  breakatwhitespace=false,     % sets if automatic breaks should only happen at 
  %whitespace
  breaklines=true,             % sets automatic line breaking
  captionpos=t,                % sets the caption-position to top
  commentstyle=\color{mygreen},   % comment style
  frame=none,               % adds a frame around the code
  keepspaces=true,            % keeps spaces in text, useful for keeping 
  %indentation of code (possibly needs columns=flexible)
  keywordstyle=\bfseries,     % keyword style
  numbers=none,               % where to put the line-numbers; possible values 
  %are (none, left, right)
  numbersep=5pt,              % how far the line-numbers are from the code
  numberstyle=\color{mygray}, % the style that is used for the line-numbers
  rulecolor=\color{black},    % if not set, the frame-color may be changed on 
  %line-breaks within not-black text (e.g. comments (green here))
  showspaces=false,           % show spaces everywhere adding particular 
  %underscores; it overrides 'showstringspaces'
  showstringspaces=false,     % underline spaces within strings only
  showtabs=false,             % show tabs within strings adding particular 
  %underscores
  stepnumber=5,               % the step between two line-numbers. If it's 1, 
  %each line will be numbered
  stringstyle=\color{red},    % string literal style
  tabsize=4,                  % sets default tabsize
  firstnumber=1      % visualizza i numeri dalla prima linea
}

% Permetti di utilizzare il grassetto per i caratteri Typewriter (per es. il 
%font di \code{...} e \file{...})
\usepackage[T1]{fontenc}
\usepackage{lmodern}


%package added
\usepackage{amsmath}

\usepackage{amsfonts}

%%%%%%%%%%%%%%
%  COSTANTI  %
%%%%%%%%%%%%%%

% In questa prima parte vanno definite le 'costanti' utilizzate soltanto da questo documento.
% Devono iniziare con una lettera maiuscola per distinguersi dalle funzioni.

\newcommand{\DocTitolo}{Relazione progetto TecWeb}

\newcommand{\DocRedazione}{I Tesori di Squitty}

\newcommand{\DocAnno}{Anno 2017/2018}

\title{\textbf{Relazione progetto TecWeb}}
\author{Roditori Pungenti}

\date{02 Luglio 2018}

\begin{document}

%\maketitle

\makeFrontPage

\tableofcontents

\newpage

\section{Abstract}

L’azienda dolciaria “I TESORI DI SQUITTY” possiede un negozio situato nel centro cittadino e un laboratorio industriale in periferia.\\
Lo scopo del progetto è quello di progettare e sviluppare un sito web che funga da vetrina e di erogare un servizio automatizzato di ordinazioni.\\

Qui di seguito sono elencati tutti i livelli di accessibilità possibili per il sito:
\begin{itemize}
	\item \textbf{Unauthenticated user}: identifica qualsiasi utente;
	\item \textbf{Authenticated user}: identifica un utente con account "al minuto"/"all'ingrosso"/"servizio" che ha effettuato l'accesso al portale web;
	\item \textbf{Authenticated administrator}: identifica un utente con account impiegato che ha effettuato l'accesso al portale web.
\end{itemize}

L'\textbf{unauthenticated user} può accedere alle pagine \textit{Home}, \textit{Per la tua casa}, \textit{Per il tuo ristorante}, \textit{Catering ed Eventi}, \textit{Accedi}, \textit{Registrati}:
\begin{itemize}
	\item \textbf{Home}: in questa pagina può trovare la storia dell'azienda dolciaria, nonchè numerose informazioni sia sul punto vendita che sullo stabilimento produttivo.

	\item \textbf{Per la tua casa}: in questa pagina può visualizzare l'immagine, gli ingredienti e la descrizione di tutti i prodotti disponibili per il commercio al dettaglio.

	\item \textbf{Per il tuo ristorante}: in questa pagina può consultare il catalogo,con relativa descrizione, dei prodotti disponibili per il commercio all'ingrosso.

	\item \textbf{Catering ed Eventi}: in questa pagina può visualizzare i servizi erogati dall'azienda.

	\item \textbf{Registrati}: in questa pagina l'utente non autenticato può effettuare la registrazione, scegliendo fra le tipologie utente: al minuto, all'ingrosso, servizio.

	\item \textbf{Accedi}: in questa pagina l'utente può autenticarsi, previo inserimento delle credenziali.\\
	Egli verrà reindirizzato alla propria area privata.
\end{itemize}

L'\textbf{authenticated user} può accedere alle sua area personale, suddivisa in \textit{Prenotazione}, \textit{Storia degli ordini}, \textit{Prodotti}.

\begin{itemize}
	\item \textbf{Prenotazione}: in questa pagina può prenotare un prodotto/servizio previo inserimento dei dati necessari, quali: quantità, data e ora di ritiro ecc...

	\item \textbf{Storia degli ordini}: in questa pagina l'utente autenticato può consultare lo storico delle prenotazioni effettuate, visualizzarne i dettagli e annallarle.

	\item \textbf{Prodotti}: in questa pagina può visualizzare il catalogo dei prodotti ordinabili compatibili con la propria tipologia di account.
\end{itemize}

Gli \textbf{authenticated administrator}, con account di tipo "Impiegato", hanno accesso alle sezioni: \textit{Ordini},\textit{Utenti}, \textit{Prodotti}.

\begin{itemize}
	\item \textbf{Ordini:} in questa pagina l'impiegato può consultare le ordinazioni divise per tipologia, visualizzarne i dettagli, cambiarne lo stato (in lavorazione ecc..) ed eliminarle.

	\item \textbf{Utenti:} in questa pagina l'impiegato visualizza l'elenco di tutti gli utenti registrati al portale web, compresi i dati anagrafici e le credenziali d'accesso e procedere alla loro eliminazione.

	\item \textbf{Prodotti:} in questa pagina l'impiegato può visualizzare i dettagli di tutti i prodotti e i servizi offerti da I Tesori di Squitty ed eliminarli.
\end{itemize}

\subsection{Analisi}

\subsubsection{Utenti destinatari}

Si sono d'apprima individuati gli utenti destinatari, tramite i seguenti step:

\begin{itemize}
	\item \textbf{Individuazione delle categorie di utenti}

	\item \textbf{Definizione di gruppo omogeneo di utenti}

	\item \textbf{Individuazione degli obiettivi per classe d'utenza}
\end{itemize}

%inserire immagini ...

\subsubsection{Base informativa}

Si è poi stilata una base informativa da cui attingere, andando a definire:

\begin{itemize}
	\item logo dell'azienda 
	\item storia dell'azienda 
	\item descrizione dello stabilimento 
	\item descrizione del negozio
	\item i prodotti
	\item i servizi
	\item contatti
\end{itemize} 

\subsection{Progettazione}

\subsection{Validazione}

Gli strumenti utilizzati per la validazione sono: \textbf{W3C Markup Validation Service} (per validazione HTML), \textbf{W3C CSS Validator} (per validazione CSS), \textbf{Total Validator} (per validazione accessibilità).
Tutte le pagine sono state validate con successo, essendo prive anche di warning.

%immagini badge validatori ...



\section{Struttura del progetto}

Per gestire la complessità del progetto e mantenere la separazione fra contenuto,stile e comportamento i file che compongono il portale sono stati raggruppati e suddivisi in 4 cartelle diverse:

\begin{itemize}
	\item \textbf{css:} questa cartella contiene tutti i file riguardanti il css, opportunamente divisi in tre file, per desktop, mobile e stampa;
	\item \textbf{php:} questa cartella contiene i file php del sito web;
	\item \textbf{img:} questa cartella contiene tutte le immagini, i loghi e le icone del sito web;
	\item \textbf{js:}  questa pagina contiene tutti gli script javascript del progetto.
\end{itemize}

\section{Gestione dei dati}

Il sito web si appoggia su una base di dati di tipo MySQL. Tutte le interazioni vengono lanciate da codice PHP tramite la libreria mysqli.\\
La struttura del database è la seguente:

%inserisci immagine ERTecWeb.jpeg

\section{Accessibilità}

La corretta sequenzialità di lettura è facilitata dall'introduzione di elementi tabindex. La presenza di colori a contrasto elevato ne facilita la lettura anche da persone con capacità visive limitate. Il sito web inoltre implementa un layout responsive, nel completo rispetto dei vigenti standard di accessibilità.\\
Nel sito web sono presenti aiuti alla navigazione, come i pulsanti "torna su", qualora fosse necessario. Non sono inoltre presenti contenuti a rischio di epilessia o che possano indurre l'utente in uno stato confusionale.\\
Abbiamo inoltre limitato l'utilizzo di Javascript soltanto alla gestione dei coockies. Questa scelta è stata fatta perseguendo la non intrusività dello scripting.\\
Tutte le pagine sono in italiano, le poche parole in inglese sono state inserite in uno <span></span> con attributo lang=”eng”.

Sono state adottate anche le seguenti misure:
\begin{itemize}
	\item Link per andare direttamente al contenuto all’interno di una access bar;
	\item Tab index regolato in modo da poter accedere ai contenuti in ordine coerente;
	\item Freccia “Torna su” fissata a fondo pagina;
	\item Presenti i meta tag per facilitare l’acesso ai motori di ricerca e fornire in modo sintetiche le info sul portale web;
	\item Link alla pagina corrente evidenziato e non cliccabile nella barra di navigazione;
	\item Link cliccabili in tutta l’area del contenitore nel quale sono contenuti e non solo nel link stesso;
	\item Sitemap gerarchica di tutte le pagine;
	\item Percorso seguito per raggiungere la pagina corrente;
	\item Contenuti dai colori contrastanti rispetto ai loro contenitori,agevolare la lettura;
	\item Link di navigazione interna alle pagine;
	\item Link già visitati di colorazione diversa rispetto a quelli non ancora visitati;
	\item Immagini chiaramente separate dal testo;
	\item Utilizzo del tool HTML5 shiv: esso permette alle versioni di IE < 9 di riconoscere i tag html5 e permetterne lo styling.
\end{itemize}


\section{Compatibilità browser}

Il presente sito web è stato testato sui principali browser. Dai test effettuati, non sono emerse particolari problematiche che possano influire negativamente sulla navigazione o sull'usabilità del sito.
Di seguito si riporta un breve riepilogo delle principali piattaforme testate:
	\begin{itemize}
		\item \textbf{Edge}: il sito viene visualizzato ed è interagibile correttamente;
		\item \textbf{Firefox}: il sito viene visualizzato ed è interagibile correttamente;
		\item \textbf{Chrome}: il sito viene visualizzato ed è interagibile correttamente;
	\end{itemize}

\section{Organizzazione del gruppo di lavoro}

La seguente tabella indica quali sono state le principali mansioni di ogni componente del gruppo. Ogni componente,tuttavia , ha contribuito almeno in parte ad ogni fase di sviluppo del progetto.

\subsection{Alessio Gobbo}

\begin{itemize}
	\item HTML
	\item PHP 
	\item CSS Desktop, Mobile
	\item stesura relazione
\end{itemize}

\subsection{Simone Ballarin}

\begin{itemize}
	\item PHP 
	\item CSS Desktop, Mobile
	\item HTML
	\item Validazione
\end{itemize}

\subsection{Gerta Lleshi}

\begin{itemize}
	\item PHP
	\item Database
	\item Javascript
	\item HTML
\end{itemize}

\subsection{Riccardo Dario}

\begin{itemize}
	\item CSS Stampa
	\item Accessibilità
	\item Stesura relazione
\end{itemize}

\end{document}